% ===================================================================
%  तालिका सं. 10 — समुद्र। — बंदरगाह।
%  Traduction 2026 d'apres texto/io, controlee sur texto/fr et
%  texto/en. Consigne : voir texto/hi/10-talika-01.tex.
%
%  Le francais decale sa numerotation d'un cran a partir de l'alinea 8.
%  C'est une coquille de l'imprimeur francais, conservee dans la
%  transcription ; le hindi suit l'ido, qui compte juste.
% ===================================================================

%%K t10-apar-1 apar
\VUsaut{4.0mm}
\VUtitre{15.0pt}{60}{तालिका सं. 10}
\VUsaut{3.0mm}
\VUpk{11.4pt}{40}{समुद्र। --- बंदरगाह।}
\VUsaut{3.0mm}

%%K t10-01-1 p
1. --- गरमियों में जब कुछ लोग पहाड़ों में शांति और ठंडी हवा खोजते हैं, दूसरे
समुद्र-तट पर जाना अधिक पसंद करते हैं। पिछले साल हमने यही किया, क्योंकि हमें
\VUgras{महासागर} \textsuperscript{(1)} बहुत भाता है।

%%K t10-02-1 p
2. --- मुझे तो \VUgras{क्षितिज} \textsuperscript{(2)} तक फैली उस विशाल जल-चादर को
निहारने में बड़ा सुख मिलता है, और यह देखने में कि विशाल \VUgras{भाप-जहाज़}
\textsuperscript{(3)} अमेरिका जाते यात्रियों को लेकर लंबी \VUgras{समुद्र-यात्रा}
पर निकलते हैं।

%%K t10-03-1 p
3. --- सो मेरे पिता, जो जानते हैं कि मुझे क्या भाता है, छुट्टियों के लिए सदा
हमारे किसी बड़े \VUgras{नौसैनिक बंदरगाह} के पास का कोई बहुत शांत तट चुनते हैं।

%%K t10-03-2 p
हमारे पिछले प्रवास में \VUgras{बेड़ा} आकर उस विशाल \VUgras{लहर-रोक}
\textsuperscript{(4)} के पीछे लंगर डाल गया जो \VUgras{नौसैनिक बंदरगाह}
\textsuperscript{(5)} को बंद करता और उसकी रक्षा करता है, और मैं इत्मीनान से
अलग-अलग \VUgras{युद्धपोत} निहार सका — छोटी \VUgras{पनडुब्बी}
\textsuperscript{(10)} से, जो डुबकी लगाकर पानी के नीचे लुप्त हो जाती है, विशाल
\VUgras{युद्धपोत} \textsuperscript{(9)} तक, जिसे अब तक \VUgras{टारपीडो नौका}
\textsuperscript{(8)} के हमलों के सिवा कुछ डराता न था।

%%K t10-tit-2 sub
\VUsaut{3.0mm}
\VUpk{10.2pt}{40}{छोटी नौकाएँ।}
\VUsaut{3.0mm}

%%K t10-04-1 p
4. --- वह पहले ही \VUgras{चतुराई} से मोटे धातु-कवच की कमज़ोर जगह ढूँढ़ लेती थी
और वहाँ वह ख़तरनाक \VUgras{टारपीडो} गाड़ देती थी जिसके विस्फोट से जहाज़ डूब जाता
है; पर वह पनडुब्बी से फिर भी कम डरावनी थी, क्योंकि विनाशक उसका पीछा सहज ही कर
सकते हैं।

%%K t10-05-1 p
5. --- मैं तेज़ बख़्तरबंद \VUgras{क्रूज़र} \textsuperscript{(6)} भी देख सका, जो
सच्चे युद्धपोत जितने ही अभेद्य हैं, कई \VUgras{सैनिक-जहाज़}
\textsuperscript{(12)}, तीन-चार \VUgras{तोप-नौकाएँ} \textsuperscript{(7)}, और अंत
में एक \VUgras{फ़्रिगेट} \textsuperscript{(11)}।
\VUgras{उस बेड़े को लंगर उठाते देखना सबसे रोचक है।}

%%K t10-06-1 p
6. --- इस्पात के उन विशाल ढेरों और उन सजीले \VUgras{तिमस्तूल} या सुडौल
\VUgras{पालदार जहाज़ों} \textsuperscript{(13)} की बनावट में कितना अंतर है, जो अब
भी व्यापारी बेड़े में काम आते हैं और जिन्हें मैं \VUgras{घाट-बाँध}
\textsuperscript{(14)} पर टहलते हुए पूरे पाल खोले बंदरगाह में आते देखता था। यह सच
है कि \VUgras{हवा} विपरीत होने पर वे अपना बहुत लाभ खो देते थे। गोदी में घुसने के
लिए उन्हें सब पाल समेटने पड़ते थे, और एक \VUgras{खींचने वाली नाव}
\textsuperscript{(16)} चाहिए होती थी जो उन्हें लेने जाए और साधारण
\VUgras{मालवाही नावों} \textsuperscript{(17)} की तरह घाट तक ले आए।

%%K t10-tit-3 sub
\VUsaut{3.0mm}
\VUpk{10.2pt}{40}{व्यापारिक बंदरगाह।}
\VUsaut{3.0mm}

%%K t10-07-1 p
7. --- जिस बंदरगाह की मैं बात कर रहा हूँ वह इसलिए इतना रोचक है कि उसमें केवल
विशाल कामों से कृत्रिम रूप से बना नौसैनिक बंदरगाह ही नहीं, बल्कि एक बड़ी नदी के
\VUgras{मुहाने} पर एक अद्भुत \VUgras{व्यापारिक बंदरगाह} भी है।

%%K t10-08-1 p
8. --- दोनों गोदियों को अलग करती भूमि की पट्टी किलेबंद है और समुद्र तक फैली
\VUgras{तोपों की क़तार} तथा \VUgras{बुर्जों} से रक्षित है। \VUgras{प्राचीरों}
\textsuperscript{(15)} पर टहलने के लिए विशेष अनुमति चाहिए। मुझे वह अपने चाचा के
ज़रिए सहज ही मिल गई थी, जो \VUgras{जहाज़-निर्माणशाला} \textsuperscript{(18)} में
अभियंता हैं, और मैं विशाल छप्परों के नीचे बनते जहाज़ देखने गया।

%%K t10-09-1 p
9. --- मैंने एक युद्धपोत को पानी में उतरते भी देखा, और मुझे वह रोमांच याद है जो
सारी भीड़ ने अनुभव किया जब वह विशाल ढेर, अपने भरोसे छोड़ा गया, अपनी पालनी पर
सरकने लगा और नदी के पानी पर तैरने आ गया।

%%K t10-10-1 p
10. --- निर्माणशाला तक पहुँचने के लिए \VUgras{परेड-मैदान} \textsuperscript{(19)}
पार करना पड़ता है, जहाँ छावनी की सेना हर दिन कवायद को आती है, और लोहे का वह
\VUgras{पुलिया} \textsuperscript{(20)} पार करना पड़ता है जो परेड-मैदान को
\VUgras{शस्त्रागार} \textsuperscript{(21)} से जोड़ता है।

%%K t10-tit-4 sub
\VUsaut{3.0mm}
\VUpk{10.2pt}{40}{शस्त्रागार और गोदियाँ।}
\VUsaut{3.0mm}

%%K t10-11-1 p
11. --- अपने चाचा के साथ मैंने वह बड़ा प्रतिष्ठान देखा, जो \VUgras{तोपों}
\textsuperscript{(22)}, \VUgras{गोलों} \textsuperscript{(23)} और \VUgras{ख़ोलों} से
भरा है। मैंने वह \VUgras{लोहशाला} \textsuperscript{(24)} भी देखी जहाँ भाप-जहाज़ों
के \VUgras{बॉयलर} \textsuperscript{(25)} बनते हैं।

%%K t10-12-1 p
12. --- बहुत पास ही \VUgras{गोदियाँ} \textsuperscript{(26)} हैं — सूखी गोदियाँ —
और वे बहुत उल्लेखनीय \VUgras{तैरती गोदियाँ} \textsuperscript{(27)}, जिनमें मैं
मरम्मत होते एक जहाज़ पर पास से किसी पोत का \VUgras{पंखा} \textsuperscript{(28)},
\VUgras{पेंदी} \textsuperscript{(29)} और \VUgras{पतवार} \textsuperscript{(30)} देख
सका।

%%K t10-13-1 p
13. --- व्यापारिक बंदरगाह नौसैनिक बंदरगाह जितना ही देखने योग्य है, और मैंने कई
घंटे उन घाटों पर बिताए जहाँ नदी चढ़ने वाले \VUgras{चक्रदार भाप-जहाज़}
\textsuperscript{(31)} बँधे रहते हैं, और उस \VUgras{तिपाई-क्रेन} \textsuperscript{(32)}
को काम करते देखते हुए, जो विशाल बोझ पंखों की तरह उठा लेती है।

%%K t10-tit-5 sub
\VUsaut{4.0mm}
\VUpk{10.2pt}{40}{कर्णधार।}
\VUsaut{3.0mm}

%%K t10-14-1 p
14. --- मैंने उन \VUgras{अंतर्महासागरीय जहाज़ों} \textsuperscript{(34)} को निहारा
जो अपनी \VUgras{झंडियाँ} \textsuperscript{(35)} फहराते बंदरगाह छोड़ते हैं, जिन्हें
एक कुशल \VUgras{कर्णधार} \textsuperscript{(36)} चलाता है, जो \VUgras{बोयों}
\textsuperscript{(40)} और \VUgras{संकेत-स्तंभों} से चिह्नित \VUgras{जलमार्ग} पर
अद्भुत ढंग से चलना जानता है, और उन \VUgras{मालनौकाओं} \textsuperscript{(41)} से
बचता है जो अपने अकेले चौकोर पाल से बड़ी कठिनाई से चलाई जाती हैं। मैं
\VUgras{अगले भाग} \textsuperscript{(37)} पर दूसरे दर्जे की \VUgras{कोठरियाँ}
\textsuperscript{(39)} पहचान सका, और \VUgras{पिछले भाग} \textsuperscript{(38)} पर
पहले दर्जे के यात्रियों के कक्ष।

%%K t10-15-1 p
15. --- कभी-कभी मैं किसी \VUgras{मालवाही जहाज़} \textsuperscript{(42)} की लदाई भी
देखता, जिसमें \VUgras{गोदाम} \textsuperscript{(43)} और \VUgras{बंदरगाह-स्टेशन}
\textsuperscript{(44)} का माल ठूँसा जाता, जो \VUgras{घाट की रेल-लाइन}
\textsuperscript{(45)} की अनेक गाड़ियाँ लाती थीं।

%%K t10-tit-6 sub
\VUsaut{3.0mm}
\VUpk{10.2pt}{40}{नौका-विहार।}
\VUsaut{3.0mm}

%%K t10-16-1 p
16. --- मैं प्रायः \VUgras{तैरती गोदी} \textsuperscript{(53)} में नाव चलाता था,
जहाँ छोटी \VUgras{नावें} \textsuperscript{(46)} शरण लेती हैं, और \VUgras{डाँड़}
\textsuperscript{(47)} चलाना मेरे लिए आनंद था। मैं \VUgras{तख़्ते}
\textsuperscript{(48)} पर चढ़ जाता, किसी \VUgras{पालदार नाव}
\textsuperscript{(49)} के,
\VUgras{रस्सियों} \textsuperscript{(52)} के सहारे \VUgras{मस्तूल}
\textsuperscript{(51)} पर \VUgras{आड़ी डाँडों} \textsuperscript{(50)} तक चढ़ता,
फिर एक \VUgras{डोंगी} \textsuperscript{(54)} में अपनी सैर जारी रखता, भारी
\VUgras{मछली की नावों} \textsuperscript{(55)} के बीच, जहाँ \VUgras{जाल}
\textsuperscript{(56)} सूखने को टँगे रहते हैं।

%%K t10-17-1 p
17. --- \VUgras{घाट} \textsuperscript{(58)} पर मेरे सामने से तरह-तरह का
\VUgras{माल} \textsuperscript{(59)} गुज़रता, \VUgras{संदूकों} \textsuperscript{(60)}
या \VUgras{गट्ठरों} \textsuperscript{(61)} में बँधा। वे मालनौकाओं का
\VUgras{लदान} \textsuperscript{(63)} बनाते थे, और शक्तिशाली \VUgras{क्रेनें}
\textsuperscript{(62)} उन्हें उठाकर \VUgras{ट्रकों} \textsuperscript{(64)} पर रखती
थीं, जब तक \VUgras{ढोने वाले} \VUgras{चुंगी-घर} \textsuperscript{(65)} में उनकी
घोषणा करते और शुल्क चुकाते।

%%K t10-18-1 p
18. --- अंत में, \VUgras{घूमने वाले पुल} \textsuperscript{(66)} या
\VUgras{जलद्वार} \textsuperscript{(69)} तक पहुँचकर, उस \VUgras{खोदने वाली नाव}
\textsuperscript{(68)} के पास से गुज़रते हुए जो \VUgras{नहर} \textsuperscript{(67)}
साफ़ करती है, मैं \VUgras{संकेत-मस्तूल} \textsuperscript{(70)} के पास घाट पर उतर
जाता।

%%K t10-19-1 p
19. --- उसी घाट की दूसरी ओर वे \VUgras{छोटी नावें} \textsuperscript{(71)} लगती हैं
जो बेड़े के अफ़सरों को किनारे लाती हैं। नाविकों को ताल से डाँड़ उठाते देखना सुंदर
दृश्य है, जब \VUgras{लहर} \textsuperscript{(72)} पर उठी नाव सहज ही उतरने की
सीढ़ियों तक आ लगती है। यहीं से \VUgras{ज्वार} और \VUgras{तट} \textsuperscript{(73)}
पर \VUgras{भाटे} तथा \VUgras{ज्वार} की चाल सबसे अच्छी देखी जा सकती है।

%%K t10-tit-7 sub
\VUsaut{3.0mm}
\VUpk{10.2pt}{40}{अफ़सरों का क्लब।}
\VUsaut{3.0mm}

%%K t10-20-1 p
20. --- घर लौटने के लिए मैं \VUgras{बिजली की ट्राम} \textsuperscript{(74)} लेता,
जिसका \VUgras{ठहराव} \textsuperscript{(75)} अफ़सरों के क्लब के सामने खड़े
\VUgras{खंभे} पर है।

%%K t10-20-2 p
क्लब के \VUgras{बालाख़ाने} \textsuperscript{(76)} से, जो \VUgras{लंगरों}
\textsuperscript{(77)}, तोपों और गोलों से सजा है, मैंने प्रायः नौसैनिक अफ़सरों का
शानदार जमघट देखा: अपनी तलवारों समेत \VUgras{लेफ़्टिनेंट}
\textsuperscript{(78)}, \VUgras{तोपख़ाने} \textsuperscript{(79)} और
\VUgras{पैदल सेना} \textsuperscript{(80)} के \VUgras{कप्तान} अपनी
\VUgras{कृपाणों} समेत। उनके पीछे एक \VUgras{नाविक} \textsuperscript{(81)} खड़ा
रहता, जो उन्हें एक \VUgras{दूरबीन} लाकर देता जब वे दूर की बात पहचानना चाहते।

%%K t10-21-1 p
21. --- मैंने जवान \VUgras{उप-लेफ़्टिनेंट} \textsuperscript{(82)} भी देखे, जो
अपने \VUgras{कमांडर} \textsuperscript{(83)} या \VUgras{एडमिरल} \textsuperscript{(84)}
से आदेश लेते थे, जब तक एक \VUgras{क्वार्टरमास्टर} \textsuperscript{(85)} उन्हें
जहाज़ तक ले जाने की प्रतीक्षा करता।

%%K t10-22-1 p
22. --- उस बालाख़ाने से दृश्य शानदार है: वहाँ से सारा तट दिखाई देता है, अपने
\VUgras{तंबुओं} \textsuperscript{(86)} और \VUgras{स्नान-कुटियों}
\textsuperscript{(87)} समेत, और नहाने के समय \VUgras{नहाने वालों}
\textsuperscript{(88)} को \VUgras{जुआघर} \textsuperscript{(89)} के सामने प्रसन्न
उछलते देखा जा सकता है।

%%K t10-23-1 p
23. --- तट के सिरे पर किनारा एक छोटा चट्टानी \VUgras{प्रायद्वीप}
\textsuperscript{(91)} बनाता है, जहाँ \VUgras{लहरें} \textsuperscript{(92)} टकराकर
टूटती हैं और जिस पर एक \VUgras{प्रकाश-स्तंभ} \textsuperscript{(93)} बना है, जो
रात में नाविकों को रास्ता दिखाता है। यह प्रायद्वीप ज़मीन से केवल एक बहुत सँकरे
\VUgras{डमरूमध्य} \textsuperscript{(90)} से जुड़ा है।

%%K t10-tit-8 sub
\VUsaut{3.0mm}
\VUpk{10.2pt}{40}{तट का सामान्य दृश्य।}
\VUsaut{3.0mm}

%%K t10-24-1 p
24. --- \VUgras{कगार} \textsuperscript{(94)} आगे चलता जाता है, समुद्र से कटा हुआ,
जो उसमें से \VUgras{खाड़ियों} \textsuperscript{(95)} और \VUgras{भृगुओं} (अंतरीपों)
की एक अनोखी लड़ी तराशता है, और उसे बहुत सुरम्य बना देता है।

%%K t10-24-2 p
\VUgras{पठार} \textsuperscript{(96)} पर, जो
\VUgras{जलडमरूमध्य} \textsuperscript{(98)} पर झुका है, एक बहुत महत्त्वपूर्ण \VUgras{गढ़} \textsuperscript{(97)}
है और कुछ कोठियाँ।

%%K t10-25-1 p
25. --- वह जलडमरूमध्य मुख्य भूमि से एक बहुत ख़तरनाक \VUgras{द्वीप}
\textsuperscript{(99)} को अलग करता है, जो \VUgras{चट्टानों} \textsuperscript{(100)}
और \VUgras{भित्तियों} से घिरा है, जहाँ नावें और बड़े जहाज़ तक कभी-कभी टूट जाते
हैं। बचाव की सारी फुरती और \VUgras{बचाव-नौकाओं} के जल्दी पहुँचने के बावजूद ये
\VUgras{जहाज़-भंग} भयानक होते हैं, और विषुव की \VUgras{आँधियों} में कई जहाज़
अपने सब आदमियों और माल समेत डूब चुके हैं।

%%K t10-25-2 p
इस पड़ोस के बावजूद \VUgras{बंदरगाह} नाविकों का बहुत आना-जाना देखता है।
\VUsaut{6.0mm}
\VUfilet{22mm}
